\documentclass[11pt,a4paper]{article}
\usepackage[margin=25mm,headheight=15pt]{geometry}
\usepackage[T1]{fontenc}
\IfFileExists{libertinus.sty}{\usepackage{libertinus}}{\usepackage{lmodern}}
\usepackage{amsmath,amssymb,amsthm,mathtools}
\usepackage{microtype,booktabs,tabularx,array,enumitem}
\usepackage{xcolor,xurl,hyperref,fancyhdr}
\hypersetup{colorlinks=true,linkcolor=blue!40!black,citecolor=green!30!black,
 urlcolor=blue!40!black,pdftitle={Resonance Profiles and Sharp Endpoint Laws for Rotating Fabius--Rvachev Distributions},
 pdfauthor={Research article prepared for Vladimir Reshetnikov},
 pdfsubject={Irrational rotations, small deviations, singular Birkhoff sums, and inverse quantiles}}
\pagestyle{fancy}\fancyhf{}
\fancyhead[L]{\small Rotating Fabius--Rvachev endpoint laws}
\fancyhead[R]{\small Research article}
\fancyfoot[C]{\thepage}
\setlength{\parskip}{0.35em}
\setlength{\parindent}{1.25em}
\setlength{\emergencystretch}{2em}
\allowdisplaybreaks[2]
\setcounter{tocdepth}{2}
\setlist[enumerate]{itemsep=0.25em,topsep=0.35em,leftmargin=2em}
\newtheorem{theorem}{Theorem}[section]
\newtheorem{proposition}[theorem]{Proposition}
\newtheorem{lemma}[theorem]{Lemma}
\newtheorem{corollary}[theorem]{Corollary}
\theoremstyle{definition}
\newtheorem{definition}[theorem]{Definition}
\newtheorem{problem}[theorem]{Research question}
\newtheorem{example}[theorem]{Example}
\theoremstyle{remark}\newtheorem{remark}[theorem]{Remark}
\newcommand{\R}{\mathbb R}
\newcommand{\N}{\mathbb N}
\newcommand{\T}{\mathbb R/\mathbb Z}
\newcommand{\PP}{\mathbb P}
\newcommand{\EE}{\mathbb E}
\newcommand{\Var}{\operatorname{Var}}
\newcommand{\Unif}{\operatorname{Unif}}
\newcommand{\dn}[1]{\left\|#1\right\|_{\T}}
\newcommand{\pos}[1]{\left(#1\right)_{+}}
\newcommand{\LL}{\mathcal L}
\newcommand{\VV}{\mathcal V}
\newcommand{\dd}{\,\mathrm d}
\newcommand{\e}{\mathrm e}
\newcommand{\ind}{\mathbf 1}
\newcommand{\bigO}{\mathcal O}
\newcommand{\repo}{\texttt{ProveIt}}
\newcommand{\at}{\alpha,\lambda}
\newcommand{\eps}{\varepsilon}
\newcommand{\pinned}{\texttt{fbba58593dc0622aa914972896150d4848f935b5}}

\title{\textbf{Resonance Profiles and Sharp Endpoint Laws\\
for Rotating Fabius--Rvachev Distributions}\\[0.55em]
\large Uniform irrational asymptotics, phase-dependent quantiles,\\
and superlinear Liouville corrections}
\author{Research article prepared for Vladimir Reshetnikov\\
\small Mathematical development and computational checks by ChatGPT}
\date{28 September 2026}

\begin{document}
\maketitle
\begin{abstract}
We investigate two explicit endpoint conjectures in the matrix-dilated
Fabius--Rvachev report in the \repo\ repository.  For independent uniform
random variables, consider
\[
 S_\beta=\sum_{n\ge1}2\e^{-\lambda n}
       |\cos\pi(n\alpha+\beta)|V_n,
 \qquad V_n\sim\Unif[0,1],\quad \lambda>0.
\]
Every irrational rotation has the same leading logarithmic Laplace and
small-deviation laws, uniformly in the phase; no Diophantine hypothesis is
needed for this leading term.  For a badly approximable rotation, we obtain
an explicit phase-uniform refinement
\[
 -\log\EE\e^{-\e^tS_\beta}
 =\frac{t^2}{2\lambda}-\frac t2-R_\beta(t)
   +O_{\alpha,\lambda}\!\bigl(\log^2(2+t)\bigr),
\]
where $R_\beta$ is a nonnegative, nondecreasing, $1$-Lipschitz profile formed
from the deepest truncated return to a cosine zero.  This proves a stronger
version of the displayed Diophantine estimate in the repository, including
phases with an exact zero.  A tilted gamma approximation supplies a relative
saddle-point formula, a three-term logarithmic cap-probability expansion,
and an explicit multiplicative inverse-quantile asymptotic.

The phase profile cannot in general be replaced by a bounded or logarithmic
carrier merely by excluding exact zeros.  Almost every phase has a common
quantile asymptotic, whereas a dense $G_\delta$ set of nonexact phases, of
Hausdorff dimension zero, has a quantile ratio with lower limit
$q=\e^{-\lambda}$ and upper limit $1$.  More generally, the limits are
classified by an inhomogeneous exponential approximation exponent.  Finally,
for a dense $G_\delta$ set of Liouville rotations we construct, for every
integer $k\ge2$, Laplace-correction subsequences of order
$t^{2-1/k}$.  These resolve the repository's Liouville-spike question in a
stronger, baseline-corrected form.  All arguments are given in ordinary
mathematical detail; numerical checks are supplementary, not proofs.
\end{abstract}
\vspace{0.3em}
\noindent\textbf{Keywords.} Fabius function; small deviations; irrational rotation;
Sudler products; exponential tilting; inverse quantiles; Liouville numbers.

\clearpage
\tableofcontents
\clearpage

\section{Research target, contribution boundary, and source status}
\label{sec:scope}

\subsection{The concrete questions being addressed}

The repository report \emph{Matrix-Dilated Fabius--Rvachev Laws}
\cite{RepoMatrix} constructs smooth laws on self-affine zonotopes and reduces
their directional boundary probabilities to infinite products of uniform
Laplace transforms.  In its section on directional caps, the source labels
\texttt{conj:diophantine-endpoint} and \texttt{conj:liouville} ask, respectively,
for control of irrational-rotation endpoint corrections and for exceptional
Liouville spikes.  The first displayed conjectural estimate, in the source's
notation, is
\begin{equation}
 \log\LL_w(s)=-\frac{(\log s)^2}{2\lambda}
       +O(\log s\,\log\log s).
 \label{eq:repoestimate}
\end{equation}
It assumes bounded continued-fraction type and a phase described as
nonresonant.  Its accompanying discussion suggests a bounded or
logarithmically growing quasiperiodic correction.  The second question asks
whether sufficiently close Liouville returns force unbounded corrections.

We resolve the displayed estimate with a stronger $O(\log s)$ error, uniformly
in \emph{all} phases, and identify the additional term responsible for that
error.  We also prove that excluding exact zeros alone does not justify the
suggested small carrier: exceptionally well approximated phases already give
linear corrections for a fixed badly approximable rotation.  The Liouville
construction gives still larger, but subquadratic, corrections.  We do
\emph{not} assert a full all-orders endpoint expansion or a bounded
quasiperiodic residual under stronger phase hypotheses.

\subsection{Relation to existing mathematics}

The scalar infinitely differentiable compactly supported function and its
probabilistic structure are classical; Arias de Reyna's exposition
\cite{Arias} is a primary reference relevant to the repository's scalar
normalization.  The products generated by
$\log|2\sin\pi(n\alpha+\gamma)|$ belong to the established theory of Sudler
products and singular Birkhoff sums.  Verschueren and Mestel
\cite{VerschuerenMestel} prove convergence along a Fibonacci subsequence and
power-law bounds for the golden rotation.  Borda \cite{Borda} studies value
distributions and concentration of Sudler products, together with Birkhoff
sums for irrational rotations.  Hauke \cite{Hauke} investigates how badly
approximable continued fractions affect extreme Sudler-product behavior.
These are related antecedents, not references establishing the present
cap-distribution and inverse-quantile conclusions.

Our use of bounded-variation discrepancy, exponential tilting, gamma
asymptotics, Borel--Cantelli, and Baire category is standard methodology.
In particular, a useful trimmed $O(\log^2 N)$ rotation estimate below is
proved directly; it is not presented as a new general theorem about Sudler
products.  The research contribution developed here is the explicit
phase-uniform endpoint profile, its transfer through a relative saddle
formula to inverse quantiles, the resulting phase classification, and the
rational-shadowing construction for the source's Liouville question.

\subsection{What the provenance check does and does not establish}

The source conjectures were inspected at repository commit
\begin{center}\small\pinned.\end{center}
The source-file blob is
\texttt{143f6da778dfc0d5ab04133005ba6715a3cf155b}.
Appendix~\ref{app:provenance} records the path and exact source labels.
The audit was targeted: it inspected the matrix report, its relevant
surrounding definitions, the repository manifest, and searches for the
conjecture label and proposed profile terminology.  It was not a proof-by-proof
review of the entire repository.  The search index and pinned-file revision
were not identical, so an unsuccessful phrase search is not treated as a
certificate of repository-wide novelty.

All theorems in this article are supported by the proofs below rather than by
assuming that the repository's other research claims are correct.  The cap
product itself is rederived.  The proofs have not been checked by a proof
assistant or independent referee.  The targeted literature check establishes
context, not universal publication priority.  Accordingly, ``proved here''
means a mathematical argument is supplied here; it does not mean that a
priority claim has been independently certified.

\begin{center}
\small
\begin{tabularx}{\textwidth}{@{}>{\raggedright\arraybackslash}p{0.25\textwidth}>{\raggedright\arraybackslash}X@{}}
\toprule
Source question or issue & Result and remaining boundary \\
\midrule
Leading irrational endpoint scale & Theorem~\ref{thm:universal}: no bounded-type
assumption is needed; convergence is phase-uniform. \\
Displayed Diophantine estimate & Theorem~\ref{thm:profile}: an explicit profile
and $O(\log^2\log s)$ residual, with total error $O(\log s)$. \\
Boundary probabilities and inversion & Theorems~\ref{thm:saddle}--\ref{thm:quantile}:
relative saddle formula and phase-corrected inverse asymptotic. \\
Meaning of ``nonresonant'' & Theorems~\ref{thm:spectrum}--\ref{thm:category}:
no exact zero is too weak for a bounded or logarithmic correction. \\
Liouville spikes & Theorem~\ref{thm:liouville}: corrections of order
$(\log s)^{2-1/k}$ along subsequences, for every integer $k\ge2$. \\
All-orders quasiperiodic expansion & Not proved; precise remaining questions
are given in Section~\ref{sec:future}. \\
\bottomrule
\end{tabularx}
\end{center}

\section{Rotating caps and a fixed normalization}
\label{sec:normalization}

Fix $\lambda>0$, set $q=\e^{-\lambda}$, and let $\alpha\in\R$ and
$\beta\in\T$.  Unless a rational comparison is explicitly under discussion,
$\alpha$ is irrational.  Let $(V_n)_{n\ge1}$ be independent random variables
with distribution $\Unif[0,1]$.  Define
\begin{align}
 a_n(\beta)&=2|\cos\pi(n\alpha+\beta)|, &
 b_n(\beta)&=q^n a_n(\beta),\label{eq:weights}\\
 S_\beta&=\sum_{n\ge1}b_n(\beta)V_n, &
 C_\beta(x)&=\PP(S_\beta\le x),\label{eq:S}\\
 \LL_\beta(s)&=\EE\e^{-sS_\beta}, &
 H_\beta(t)&=-\log\LL_\beta(\e^t).
 \label{eq:H}
\end{align}
The series converges for every sample point, since
$\sum b_n\le2q/(1-q)$.  All logarithms are natural.

\begin{proposition}[Reduction of a directional cap]
\label{prop:cap}
Let $U_n$ be independent $\Unif[-1,1]$ variables, let $R_\theta$ denote planar
rotation through $\theta$, and let
\[
 X=\sum_{n\ge1}q^nR_{n\theta}v\,U_n,
 \qquad v\in\R^2\setminus\{0\}.
\]
For $w\ne0$, put
$h(w)=\sum_{n\ge1}q^n|\langle w,R_{n\theta}v\rangle|$.
If $\alpha=\theta/\pi$ and $\delta$ is the angle of $v$ minus the angle of
$w$, then
\[
 h(w)-\langle w,X\rangle\ \overset{d}=\
 \|v\|\|w\|S_\beta,\qquad \beta=\delta/\pi\pmod1.
\]
Moreover,
\begin{equation}
 \LL_\beta(s)=\prod_{n\ge1:b_n>0}
            \frac{1-\e^{-sb_n}}{sb_n},\qquad s\ge0.
 \label{eq:product}
\end{equation}
A zero coefficient contributes a factor $1$.
\end{proposition}
\begin{proof}
The support bound in direction $w$ is the sum of the absolute projected
segment lengths, and every finite choice of segment endpoints is approximable
within the support, so its supremum is $h(w)$.  For a real coefficient $c$,
$|c|-cU$ has the same law as $2|c|V$.  Apply this independently to every
projected term and use
$\langle w,R_{n\theta}v\rangle=\|v\|\|w\|\cos(n\theta+\delta)$.
The finite-product formula follows by integration of a uniform variable.
Bounded convergence of the exponential as the finite sum tends to $S_\beta$
gives the infinite-product identity.
\end{proof}

The factor $\|v\|\|w\|$ is not silently absorbed into $\lambda$ or $\beta$.
For a physical deficit $D=A S_\beta$ with $A>0$, substitute $s\mapsto As$ in
the Laplace formulas and $x\mapsto x/A$ in the probability formulas.  Its
quantile is exactly $A C_\beta^{-1}(y)$.

Write $\dn{u}=\min_{m\in\mathbb Z}|u-m|$ and
$d_n(\beta)=\dn{n\alpha+\beta-1/2}$.  For $0\le d\le1/2$,
\begin{equation}
 4d\le2\sin(\pi d)\le2\pi d.
 \label{eq:distance-amplitude}
\end{equation}
A phase is called \emph{exact} if $d_m(\beta)=0$ for some $m\ge1$, and
\emph{nonexact} otherwise.  Irrationality permits at most one exact index.
This terminology deliberately avoids imposing an unstated quantitative
meaning on the word ``nonresonant.''

\begin{lemma}[Support and quantiles]
\label{lem:quantile-exists}
For irrational $\alpha$, $S_\beta$ has no atoms and its support is the whole
interval $[0,\sum b_n]$.  Its distribution function is continuous and strictly
increasing on that interval.  Thus $C_\beta^{-1}(y)$ is unambiguous for
$0<y<1$.
\end{lemma}
\begin{proof}
At least one coefficient is positive, so convolution with its uniform law
eliminates atoms.  Finite sums have full interval supports; their deterministic
tail bound tends to zero.  More explicitly, any interior point can be
approximated by a finite sum with values in the interiors of its positive
uniform intervals, while the remaining tail is made smaller than a prescribed
neighborhood.  This gives positive probability to every open subinterval of
the claimed support.  The stated strict increase follows.
\end{proof}

\section{The soft cutoff and the scalar baseline}
\label{sec:kernel}

Set
\begin{equation}
 F(u)=\log\frac{u}{1-\e^{-u}}\quad(u>0),\quad F(0)=0,
 \qquad h(z)=F(\e^z),\quad h(-\infty)=0.
 \label{eq:kernel}
\end{equation}
The infinite product becomes
\begin{equation}
 H_\beta(t)=\sum_{n\ge1}h(t-\lambda n+f_n),
 \qquad f_n=\log a_n\in[-\infty,\log2].
 \label{eq:Hsum}
\end{equation}

\begin{lemma}[Elementary kernel estimates]
\label{lem:kernel}
There is an absolute constant $C$ such that
\begin{align}
 &0\le F(u)\le u/2,\qquad 0\le F'(u)\le1/2,\label{eq:Fbounds}\\
 &0\le h(z)-z_+\le C\e^{-|z|}\qquad(z\in\R),\label{eq:hcut}\\
 &h'(z)=J(\e^z),\qquad
 J(u)=1-\frac{u}{\e^u-1},\label{eq:J}\\
 &J(u)-uJ'(u)=V(u),\qquad
 V(u)=1-\frac{u^2\e^u}{(\e^u-1)^2}.
 \label{eq:V}
\end{align}
At $u=0$, $J(0)=V(0)=0$.  For $u\ge0$,
$0\le J(u)\le1$, $0\le V(u)\le1$, and
\[
 J(u)=O(u),\quad V(u)=O(u^2)\quad(u\downarrow0),
 \qquad J(u),V(u)\longrightarrow1\quad(u\to\infty)
\]
with exponentially small errors at infinity, up to polynomial factors.
\end{lemma}
\begin{proof}
$F(u)$ is the negative logarithm of the Laplace transform of $\Unif[0,1]$.
Jensen's inequality gives $F(u)\le u/2$.  Its derivative is the mean of this
uniform variable under exponential tilt; it lies between $0$ and its untilted
mean $1/2$, since the derivative of that mean is the negative tilted variance.
For $z\le0$, $h(z)\le\e^z/2$.  For $z\ge0$,
\[
 h(z)-z=-\log(1-\e^{-\e^z})\ge0;
\]
this is bounded by $C\e^{-z}$, with a larger constant covering $z=0$.
Differentiation gives \eqref{eq:J} and \eqref{eq:V}.

Under tilt, $uV_1$ has density
$\e^{-w}/(1-\e^{-u})$ on $[0,u]$.  Its mean and variance are $J(u)$ and
$V(u)$, respectively.  This proves nonnegativity.  The displayed formulas
prove the upper bounds by $1$; alternatively the mean is at most the mean of
an unconditioned unit exponential.  Taylor expansion at zero and the explicit
formulas at infinity prove the remaining estimates.
\end{proof}

Define, for $t\ge0$,
\begin{equation}
 B_\lambda(t)=\sum_{n\ge1}(t-\lambda n)_+.
 \label{eq:Bdef}
\end{equation}
Writing $N=\lfloor t/\lambda\rfloor$ and $\eta=\{t/\lambda\}$ gives the
\emph{exact} identity
\begin{equation}
 B_\lambda(t)=Nt-\frac{\lambda N(N+1)}2
 =\frac{t^2}{2\lambda}-\frac t2+
       \frac\lambda2\eta(1-\eta).
 \label{eq:Bexact}
\end{equation}
Thus the linear term $-t/2$ is forced by the indexing $n\ge1$; it is not a
resonance effect.

For later use, the kernel bounds imply
\begin{equation}
 \sum_{n\ge1}h(t-an)=B_a(t)+O_\lambda(1)
 \quad\text{uniformly for }a\ge\lambda,
 \label{eq:scalarcut}
\end{equation}
since the exponentially decaying errors in \eqref{eq:hcut} can be summed over
any lattice with spacing at least $\lambda$.  They also give the useful
truncation certificate
\begin{equation}
 0\le H_\beta(t)-\sum_{n=1}^{M}F(\e^t b_n)
 \le\frac{\e^tq^{M+1}}{1-q}.
 \label{eq:tail}
\end{equation}
The same bound holds for the omitted part of $H_\beta'(t)$, because
$J(u)=uF'(u)\le u/2$.  A safe variance-tail bound is
$4\e^{2t}q^{2(M+1)}/(1-q^2)$, using $V(u)\le u^2$.

\section{The universal leading law for every irrational rotation}
\label{sec:universal}

\begin{theorem}[Phase-uniform irrational universality]
\label{thm:universal}
For every fixed irrational $\alpha$ and every $\lambda>0$, uniformly in
$\beta\in\T$,
\begin{align}
 H_\beta(t)&=\frac{t^2}{2\lambda}+o(t^2),\qquad t\to\infty,
 \label{eq:universalH}\\
 \log C_\beta(x)&=-\frac{\log^2(1/x)}{2\lambda}
                  +o\!\bigl(\log^2(1/x)\bigr),\qquad x\downarrow0,
 \label{eq:universalC}\\
 -\log C_\beta^{-1}(y)&\sim\sqrt{2\lambda\log(1/y)},\qquad y\downarrow0.
 \label{eq:universalQ}
\end{align}
The little-$o$ terms need not be uniform over irrational $\alpha$.
\end{theorem}
\begin{proof}
Since $a_n\le2$, monotonicity of $h$ and \eqref{eq:scalarcut} imply
\[
 H_\beta(t)\le B_\lambda(t+\log2)+O_\lambda(1)
             =\frac{t^2}{2\lambda}+O_\lambda(t).
\]
For the lower bound, fix $0<\eps<1$.  The set
$E_\eps=\{u\in\T:2|\cos\pi u|<\eps\}$ is an interval about $1/2$,
of length tending to zero as $\eps\downarrow0$.  Irrational rotation gives
\begin{equation}
 \#\{1\le n\le N:n\alpha+\beta\in E_\eps\}
 =|E_\eps|N+o(N)
 \label{eq:uniformequi}
\end{equation}
uniformly in $\beta$.  To see the uniformity directly, a nonconstant Fourier
mode has average bounded by
$2/(N|1-\e^{2\pi i k\alpha}|)$, independently of the starting phase.
Approximate the interval indicator from above and below by continuous
periodic functions and then by finite trigonometric polynomials.

Retain the terms with $a_n\ge\eps$, and put
$N=\lfloor(t+\log\eps)/\lambda\rfloor$.  The inequality $h(z)\ge z_+$
gives
\[
 H_\beta(t)\ge B_\lambda(t+\log\eps)
  -(t+|\log\eps|)\bigl(|E_\eps|N+o(N)\bigr).
\]
Divide by $t^2$, first let $t\to\infty$, and then let $\eps\downarrow0$.
This proves \eqref{eq:universalH} uniformly in phase.

Set $L=\log(1/x)$.  Chernoff's inequality at $s=1/x$ gives
\begin{equation}
 C_\beta(x)\le\exp(1-H_\beta(L)).
 \label{eq:chernoff}
\end{equation}
For a phase-independent lower bound, couple
$S_\beta\le S^*=2\sum q^nV_n$.  Choose
\[
 N=\left\lceil\frac{L+\log(4/(1-q))}{\lambda}\right\rceil.
\]
The deterministic tail of $S^*$ after $N$ is at most $x/2$.
For sufficiently large $L$, each threshold $x/(4Nq^n)$, $1\le n\le N$,
lies in $(0,1)$.  On the event $V_n\le x/(4Nq^n)$ for all these indices,
the first $N$ terms contribute at most $x/2$.  Independence therefore gives
\begin{align*}
 \log C_\beta(x)
 &\ge -NL-N\log(4N)+\frac{\lambda N(N+1)}2\\
 &=-\frac{L^2}{2\lambda}-\frac{L\log L}{\lambda}+O_\lambda(L).
\end{align*}
Together with \eqref{eq:chernoff}, this proves \eqref{eq:universalC}.
Monotone inversion of the two uniform quadratic bounds proves
\eqref{eq:universalQ}; Lemma~\ref{lem:quantile-exists} supplies the inverse.
\end{proof}

\begin{proposition}[Rational comparison]
\label{prop:rationalleading}
Let $\alpha=p/r$ be reduced, and let $\nu$ of the $r$ periodic amplitudes
$a_1,\ldots,a_r$ be positive.  If $\nu>0$, then
\[
 H_\beta(t)\sim\frac{\nu}{r}\frac{t^2}{2\lambda},\qquad
 \log C_\beta(x)\sim-\frac{\nu}{r}\frac{\log^2(1/x)}{2\lambda}.
\]
For $r\ge2$, a generic phase has $\nu=r$, while a phase orthogonal to one
orbit direction has $\nu=r-1$.
\end{proposition}
\begin{proof}
For each positive residue $j$, its contribution is
$\sum_{k\ge0}h(t-\lambda j-\lambda rk+\log a_j)$, whose leading term is
$t^2/(2\lambda r)$ by the scalar cutoff.  Sum over the $\nu$ residues.
Chernoff gives the probability upper bound.  The finite-coordinate lower-bound
event used above now needs only the positive indices: their count through
$N$ is $\nu N/r+O_r(1)$ and their index sum is
$\nu N^2/(2r)+O_r(N)$.  The fixed nonzero amplitude factors change the
logarithm by only $O_{\alpha,\beta,\lambda}(N)$, proving the lower bound.
Distinct residues yield distinct points modulo one, so at most one amplitude
vanishes when $r\ge2$.
\end{proof}

This rational/irrational difference is relevant to the Liouville construction:
a very accurate rational approximation can temporarily suppress a positive
fraction of the effective digits, although that fraction tends to zero along
increasing rational denominators.

\section{A uniform trimmed rotation estimate}
\label{sec:rotation}

For the next five sections impose the bounded-type hypothesis
\begin{equation}
 \dn{m\alpha}\ge\frac{\kappa}{m}\qquad(m\ge1)
 \quad\text{for some fixed }\kappa>0.
 \tag{BT}\label{eq:BT}
\end{equation}
This is the usual badly approximable condition, equivalent to bounded
continued-fraction partial quotients.  Constants may depend on $\alpha$
(or a chosen bound for its partial quotients) and $\lambda$, but never on
$\beta$ unless explicitly indicated.

\begin{lemma}[Uniform interval discrepancy]
\label{lem:discrepancy}
Under \eqref{eq:BT}, there is $D_\alpha$ such that, for every interval
$I\subset\T$, $N\ge1$, and $\beta\in\T$,
\[
 \left|\#\{1\le n\le N:n\alpha+\beta\in I\}-N|I|\right|
 \le D_\alpha\log(2N).
\]
Consequently, for a periodic function $g$ of bounded variation,
\begin{equation}
 \left|\sum_{n=1}^N g(n\alpha+\beta)-N\int_\T g\right|
 \le D_\alpha\operatorname{Var}(g)\log(2N),
 \label{eq:BV}
\end{equation}
with a harmless change of the constant.
\end{lemma}
\begin{proof}
Let $p_j/Q_j$ be a continued-fraction convergent.  The first $Q_j$ points
of an orbit block differ, after a permutation, from a translated equally
spaced $Q_j$-point grid by less than $1/Q_{j+1}\le1/Q_j$ on the circle.
Indeed, $|\alpha-p_j/Q_j|<1/(Q_jQ_{j+1})$, and the accumulated error over a
block of $Q_j$ successive indices is at most $1/Q_{j+1}$ after fixing its
starting phase.  The grid count in an interval differs from $Q_j|I|$ by
at most two.  Only a bounded number of grid points can lie within $1/Q_j$
of the interval's endpoints.  Thus a block of length $Q_j$ has an absolute
counting error bounded by a universal constant, uniformly in its starting
phase.

If all partial quotients are bounded by $A$, the greedy continued-fraction
expansion of $N$ decomposes its first $N$ orbit points into at most $A+1$
blocks of each convergent length.  The denominators grow at least as fast
as the Fibonacci denominators, so there are $O(\log(2N))$ different lengths.
Adding the block errors proves the interval estimate.  Finally, apply
Stieltjes integration by parts to the signed counting measure minus
$N$ times Lebesgue measure; its total mass is zero and its interval
cumulative function is bounded by the estimate just proved.  This gives
\eqref{eq:BV}.
\end{proof}

\begin{lemma}[One exceptional term suffices]
\label{lem:trim}
Let $f(u)=\log(2|\cos\pi u|)$, with value $-\infty$ at $1/2$.
For $N\ge2$, choose $m\in\{1,\ldots,N\}$ minimizing
$d_n=\dn{n\alpha+\beta-1/2}$.  Under \eqref{eq:BT},
\begin{align}
 d_n&\ge\frac{\kappa}{2N}\quad(n\ne m),\label{eq:spacing}\\
 f(n\alpha+\beta)&\ge-\log N-C_\alpha\quad(n\ne m),
 \label{eq:fmin}\\
 \left|\sum_{\substack{1\le n\le N\\n\ne m}}
               f(n\alpha+\beta)\right|
 &\le C_\alpha\log^2(2N).
 \label{eq:trim}
\end{align}
The sum in \eqref{eq:trim} is finite even at an exact phase, since its only
possible infinite term was removed.
\end{lemma}
\begin{proof}
Any two of the first $N$ orbit points are separated by at least $\kappa/N$.
If $d_m\ge\kappa/(2N)$, the first assertion is immediate.  Otherwise the
triangle inequality and the separation bound give
$d_n\ge\kappa/N-d_m>\kappa/(2N)$ for every $n\ne m$.
Inequality~\eqref{eq:distance-amplitude} gives \eqref{eq:fmin}.

Choose a constant $C_0$ large enough that
$\tau_N=-\log N-C_0$ is no larger than the lower bound in
\eqref{eq:fmin}, and put $g_N=\max(f,\tau_N)$.
Its variation on the circle is $O_\alpha(\log(2N))$.  Also
\begin{equation}
 \int_\T f=0,\qquad
 0\le\int_\T(g_N-f)=O_\alpha(1/N).
 \label{eq:truncatedmean}
\end{equation}
For completeness, the mean-zero identity follows from
$\int_0^{\pi/2}\log\sin u\dd u=-(\pi/2)\log2$ and reflection.
The integral identity for sine is obtained by adding its sine and cosine
versions and substituting $2u$ in their product.  The truncation correction
is supported where the distance to $1/2$ is $O(1/N)$; integration of
$\log(c/(Nd))$ over such an interval is $O(1/N)$, proving the second claim.

Equation~\eqref{eq:BV} now gives
$\sum_{n\le N}g_N(n\alpha+\beta)=O_\alpha(\log^2(2N))$.
Except at $m$, $g_N$ agrees with $f$; and
$|g_N(m\alpha+\beta)|=O_\alpha(\log(2N))$.
Subtract this one finite term.  This proves \eqref{eq:trim}.
\end{proof}

The uniformity in this lemma is important.  An estimate for a fixed unshifted
Sudler product would not, without an additional argument, control a phase
chosen extremely close to one of the first $N$ singular orbit locations.

\section{The resonance profile and the sharp Laplace law}
\label{sec:profile}

\begin{definition}[Truncated resonance profile]
\label{def:profile}
Let
\begin{equation}
 k_n(\beta)=\bigl(-\log a_n(\beta)\bigr)_+\in[0,\infty],
 \qquad
 R_\beta(t)=\max_{n\ge1}\min\{k_n(\beta),(t-\lambda n)_+\}.
 \label{eq:profiledef}
\end{equation}
The conventions are $\min\{\infty,u\}=u$ for finite $u$ and $R_\beta(t)=0$
for $0\le t<\lambda$.  Only finitely many indices contribute when $t$ is fixed.
\end{definition}

Each function in the maximum grows with slope $1$ after its starting time
$\lambda n$ until it reaches height $k_n$, after which it stays constant.
Thus
\begin{equation}
 0\le R_\beta(t)\le t,\qquad
 0\le R_\beta(t_2)-R_\beta(t_1)\le t_2-t_1
 \quad(t_2\ge t_1).
 \label{eq:Lipschitz}
\end{equation}
This elementary finite maximum, rather than an untruncated singular sum,
is the correction that survives uniformly over arbitrary phases.

\begin{theorem}[Uniform profile expansion]
\label{thm:profile}
Under \eqref{eq:BT}, as $t\to\infty$,
\begin{equation}
 H_\beta(t)=B_\lambda(t)-R_\beta(t)
                   +O_{\at}\!\bigl(\log^2(2+t)\bigr)
 \label{eq:mainprofile}
\end{equation}
uniformly for every $\beta\in\T$, including exact phases.  Equivalently,
\begin{equation}
 \log\LL_\beta(\e^t)
 =-\frac{t^2}{2\lambda}+\frac t2+R_\beta(t)
                   +O_{\at}\!\bigl(\log^2(2+t)\bigr).
 \label{eq:laplaceprofile}
\end{equation}
\end{theorem}
\begin{proof}
Let $N=\lfloor t/\lambda\rfloor$ and choose the exceptional index $m$ from
Lemma~\ref{lem:trim}.  Set $a=t-\lambda m\ge0$.
First consider all $n\le N$ other than $m$.  Their $f_n$ lie between
$-\log N-C_\alpha$ and $\log2$.
Except in a final interval of $O_{\at}(\log(2N))$ indices, the quantity
$t-\lambda n+f_n$ is at least $3\log(2N)$.  On that early interval,
\eqref{eq:hcut} permits replacing $h(t-\lambda n+f_n)$ by
$t-\lambda n+f_n$ with a total error $O_{\at}(1)$.
On the final interval both the negative part of this linear expression and
the soft-cutoff error are $O_{\at}(\log(2N))$ per term.  Consequently,
\begin{align}
 \sum_{\substack{n\le N\\n\ne m}}h(t-\lambda n+f_n)
 &=\sum_{\substack{n\le N\\n\ne m}}(t-\lambda n)
   +\sum_{\substack{n\le N\\n\ne m}}f_n
   +O_{\at}(\log^2(2N))\nonumber\\
 &=B_\lambda(t)-a+O_{\at}(\log^2(2N)).
 \label{eq:trimmedH}
\end{align}
For $n>N$, $h(t-\lambda n+f_n)\le\e^{t-\lambda n}$, so the total tail is
$O_\lambda(1)$.

If $f_m\le0$, including $f_m=-\infty$, the cutoff estimate gives
\[
 h(a+f_m)-a=-\min\{k_m,a\}+O(1).
\]
For $f_m>0$, the same statement holds with $k_m=0$ and an $O(1)$ error,
since $f_m\le\log2$ and $h$ is $1$-Lipschitz.
Combining this with \eqref{eq:trimmedH} yields
\[
 H_\beta(t)=B_\lambda(t)-\min\{k_m,t-\lambda m\}
                           +O_{\at}(\log^2(2N)).
\]
All other indices $n\le N$ have $k_n\le\log N+C_\alpha$.  Thus replacing
the exceptional term by the maximum $R_\beta(t)$ changes the expression by
at most $O_\alpha(\log(2N))$.  Equation~\eqref{eq:Bexact} finishes the proof.
\end{proof}

\begin{corollary}[Resolution of the displayed source estimate]
\label{cor:repo}
Under \eqref{eq:BT}, uniformly in phase,
\[
 \log\LL_\beta(s)=-\frac{(\log s)^2}{2\lambda}+O_{\at}(\log s).
\]
In particular, the estimate \eqref{eq:repoestimate} holds, with no phase
exclusion at all.
\end{corollary}
\begin{proof}
Use $0\le R_\beta(t)\le t$ and $\log^2 t=o(t)$ in
Theorem~\ref{thm:profile}.
\end{proof}

\begin{corollary}[A quantitative nonresonance condition]
\label{cor:polynomial}
Suppose that, in addition to \eqref{eq:BT}, there are $c>0$ and $\tau>0$ with
\begin{equation}
 \dn{n\alpha+\beta-1/2}\ge c n^{-\tau}\qquad(n\ge1).
 \label{eq:polynomialphase}
\end{equation}
Then $R_\beta(t)=O_{c,\tau,\lambda}(\log(2+t))$ and
\[
 \log\LL_\beta(\e^t)=-\frac{t^2}{2\lambda}+\frac t2
                     +O_{\alpha,\lambda,c,\tau}(\log^2(2+t)).
\]
\end{corollary}
\begin{proof}
Inequality~\eqref{eq:distance-amplitude} gives
$k_n\le\tau\log n+C_c$.  In \eqref{eq:profiledef}, only
$n\le t/\lambda$ matter.
\end{proof}

\begin{remark}[Why the untruncated cocycle is not the uniform correction]
The source's cocycle is $\sum_{n\le N}f_n$.  A single $f_m$ can be arbitrarily
negative, or $-\infty$, while its actual Laplace contribution stays between
$0$ and a quantity of order $t-\lambda m$.  The minimum in
\eqref{eq:profiledef} is precisely this saturation.  Before saturation, a
small digit is effectively missing; after saturation, its full logarithmic
amplitude enters.  The formula is therefore compatible with continuity of the
Laplace transform at a phase where a digit becomes exactly zero.
\end{remark}

\section{A relative saddle formula from tilted gamma blocks}
\label{sec:saddle}

A logarithmic Laplace estimate by itself is not a relative asymptotic for a
small probability.  This section supplies the missing local argument.  For
$s=\e^t$, define the tilted probability measure
\begin{equation}
 \frac{\dd\PP_t}{\dd\PP}=\exp(-sS_\beta+H_\beta(t)).
 \label{eq:tilt}
\end{equation}
Expectations and variances under this measure are denoted by $\EE_t$ and
$\Var_t$.  Set
\begin{equation}
 m_\beta(t)=H_\beta'(t),\qquad
 \VV_\beta(t)=H_\beta'(t)-H_\beta''(t)
             =s^2\Var_t(S_\beta).
 \label{eq:tiltedparameters}
\end{equation}
The variance is strictly positive.

\begin{lemma}[Tilted summands and their effective number]
\label{lem:tilted}
Under \eqref{eq:BT}, uniformly in $\beta$ as $t\to\infty$,
\begin{equation}
 H_\beta'(t)=\frac t\lambda+O_{\at}(\log(2+t)),\qquad
 \VV_\beta(t)=\frac t\lambda+O_{\at}(\log(2+t)).
 \label{eq:derivativeasymptotic}
\end{equation}
More specifically, write $u_n=s b_n$ and $W_n=s b_nV_n$ under $\PP_t$.
The nonzero $W_n$ are independent unit exponential variables conditioned to
lie in $[0,u_n]$.  Their mean and variance are $J(u_n)$ and $V(u_n)$.
If $N=\lfloor t/\lambda\rfloor$, then all but one of the indices
\[
 1\le n\le N-C_{\at}\log(2N)
\]
have $u_n\ge N^3$, after increasing the constant and taking $N$ large.
The sum of the means and the sum of the variances of all remaining terms
are $O_{\at}(\log(2N))$.
\end{lemma}
\begin{proof}
Tilting a product of independent laws preserves independence: for any finite
set of indices, integrate the complementary factors in
\eqref{eq:tilt} and cancel their Laplace transforms.  A uniform variable on
$[0,u]$ becomes a variable with density
$\e^{-w}/(1-\e^{-u})$ on that interval.  The formulas for $J$ and $V$
therefore apply, and summing them proves \eqref{eq:tiltedparameters}.
Termwise differentiation is also justified directly by the summable small-$u$
bounds on any compact $t$ interval.

Use the same exceptional index as in Lemma~\ref{lem:trim}.  For every other
$n\le N$, $a_n\ge\e^{-C_\alpha}/N$.  If
$\lambda(N-n)\ge4\log N+C_\alpha$, then
$u_n=\e^{t-\lambda n}a_n\ge N^3$.
There are $N-O_{\at}(\log(2N))$ such indices.  For them, both mean and
variance are $1$ with exponentially small errors.  The remaining indices
through $N$ number $O_{\at}(\log(2N))$, and each has mean and variance at
most $1$.  For $n>N$, the bounds $J(u)=O(u)$ and $V(u)=O(u^2)$ give
bounded geometric tails.  Summing proves all assertions.
\end{proof}

\begin{lemma}[A gamma block gives the needed local limit]
\label{lem:gamma}
Suppose $G_M$ has the gamma density
\[
 g_M(w)=\frac{w^{M-1}\e^{-w}}{\Gamma(M)}\ind_{w>0},
 \qquad M\to\infty,
\]
and $B_M$ is independent of $G_M$ with
$\Var(B_M)=O(\log M)$.  Let
$a_M=M+\EE B_M+o(1)$.  Then
\begin{equation}
 \EE\left[\e^{G_M+B_M-a_M}
                \ind_{G_M+B_M\le a_M}\right]
 =\frac{1+o(1)}{\sqrt{2\pi M}}.
 \label{eq:gammalocal}
\end{equation}
The conclusion is uniform in any auxiliary parameters for which the stated
variance bound and $o(1)$ are uniform.  It remains true if the law of
$G_M+B_M$ is changed by total variation $o(M^{-1/2})$, provided the same
centering $a_M$ is used.
\end{lemma}
\begin{proof}
Stirling's formula implies
\begin{equation}
 g_M(M+y)=\frac{1+o(1)}{\sqrt{2\pi M}}
       \quad\text{uniformly for }|y|\le M^{1/4},
 \qquad \sup_w g_M(w)\le\frac{C}{\sqrt M}.
 \label{eq:gammastirling}
\end{equation}
For example, substitute $w=M+y$ in the logarithm of the density; the quadratic
term is $-y^2/(2M)=O(M^{-1/2})$ and the remaining terms tend to zero uniformly
on that window.  The supremum occurs at $M-1$ when $M>1$ and follows from
the same formula.

Let $\widetilde g_M$ be the density of $G_M+B_M$.  Conditioning on $B_M$,
\[
 \widetilde g_M(a_M-v)=\EE\,g_M(a_M-v-B_M).
\]
Chebyshev's inequality gives
$\PP(|B_M-\EE B_M|>M^{1/4}/3)=O(\log M/\sqrt M)=o(1)$.
For $0\le v\le M^{1/4}/3$, the complementary event puts the gamma density
inside the window in \eqref{eq:gammastirling}, for large $M$.
The exceptional event contributes $o(M^{-1/2})$ by the global density bound.
Thus $\widetilde g_M(a_M-v)=(2\pi M)^{-1/2}(1+o(1))$ uniformly on that
$v$ interval, and globally it is bounded by $C/\sqrt M$.
Finally,
\[
 \EE[\e^{G_M+B_M-a_M}\ind_{G_M+B_M\le a_M}]
 =\int_0^\infty\e^{-v}\widetilde g_M(a_M-v)\dd v.
\]
The tail after $M^{1/4}/3$ is negligible, proving the claim.
The integrand defining the expectation is bounded by $1$, so a total
variation perturbation of the stated order changes it by $o(M^{-1/2})$.
\end{proof}

\begin{theorem}[Uniform relative cap saddle]
\label{thm:saddle}
Assume \eqref{eq:BT}.  For sufficiently small $x>0$, let $t_x=t_x(\beta)$
be the unique solution of
\begin{equation}
 \e^{t_x}x=H_\beta'(t_x).
 \label{eq:saddle-equation}
\end{equation}
Then, uniformly in $\beta\in\T$ as $x\downarrow0$,
\begin{equation}
 C_\beta(x)=
 \frac{\exp\!\bigl(\e^{t_x}x-H_\beta(t_x)\bigr)}
      {\sqrt{2\pi\VV_\beta(t_x)}}\,(1+o(1)).
 \label{eq:saddle}
\end{equation}
\end{theorem}
\begin{proof}
The tilted mean $\EE_sS_\beta$ is strictly decreasing in $s$, because its
derivative is $-\Var_s(S_\beta)<0$.  It tends to zero as $s\to\infty$,
as also follows from Lemma~\ref{lem:tilted}.  Its value at $s=0$ is positive.
For irrational $\alpha$, two successive cosine factors cannot both vanish;
compactness in $\beta$ gives a positive uniform lower bound for the sum of
their absolute values.  The first two coefficients therefore give a positive
uniform lower bound for $\EE S_\beta$.  It follows that for all sufficiently
small $x$, the equation $\EE_sS_\beta=x$ has a unique solution.  That solution
has $t_x\to\infty$ uniformly: for bounded $s$, positivity and continuity in
$\beta$, again using the first two coefficients, bound the tilted mean away
from zero.  Equation~\eqref{eq:saddle-equation} is its logarithmic form.

We now prove the required local statement at an arbitrary large $t$.
Use the early indices of Lemma~\ref{lem:tilted}, excluding its possible
exceptional index.  Let their number be $M=N-O_{\at}(\log(2N))$.
Each early summand is a unit exponential conditioned to be at most $u_n$,
where $u_n\ge N^3$.  The total variation distance between that conditioned
law and the unconditioned exponential law is $\e^{-u_n}$.  A product coupling
therefore replaces the entire early sum by an independent gamma variable
$G_M$ with total variation error at most $N\e^{-N^3}$.

Let $B_M$ be the sum of all the other tilted summands.  It is independent
of the replaced early block and has mean and variance $O_{\at}(\log(2N))$.
The replacement changes the mean of the early block by at most
\[
 \sum_{\mathrm{early}\ n}\frac{u_n}{\e^{u_n}-1}
 =O(N^4\e^{-N^3})=o(1).
\]
Consequently the original dimensionless sum $W=sS_\beta$ has center
$m_\beta(t)=M+\EE B_M+o(1)$ and its law is exponentially close in total
variation to $G_M+B_M$.  Lemma~\ref{lem:gamma} yields
\begin{equation}
 \EE_t[\e^{W-m_\beta(t)}\ind_{W\le m_\beta(t)}]
 =\frac{1+o(1)}{\sqrt{2\pi\VV_\beta(t)}},
 \label{eq:localtilt}
\end{equation}
since $M\sim\VV_\beta(t)$ uniformly.  This proves more than a weak central
limit theorem: it gives the exponentially weighted, one-sided local estimate
needed for a relative probability.

At $t=t_x$, $sx=m_\beta(t_x)$.  Reversing the tilt gives the exact identity
\[
 C_\beta(x)=\e^{sx-H_\beta(t_x)}
       \EE_{t_x}[\e^{sS_\beta-sx}\ind_{sS_\beta\le sx}].
\]
Insert \eqref{eq:localtilt}.  Every preceding estimate is phase-uniform,
so the relative $o(1)$ in \eqref{eq:saddle} is phase-uniform as well.
\end{proof}

\begin{remark}[Scope of the relative formula]
Theorem~\ref{thm:saddle} retains the \emph{exact} functions $H_\beta$ and
$\VV_\beta$ at the exact saddle.  Replacing $H_\beta$ by an approximation
with $O(\log^2 t)$ error does not preserve a relative $1+o(1)$ probability
estimate.  It does preserve the logarithmic expansion below, and that
logarithmic accuracy is enough to give a relative inverse-quantile formula.
\end{remark}

\section{Three-term cap asymptotics and inverse quantiles}
\label{sec:quantile}

Define the constant
\begin{equation}
 c_\lambda=\frac{1+\log\lambda}{\lambda}+\frac12.
 \label{eq:clambda}
\end{equation}
For small $x>0$, let
\begin{equation}
 L=\log(1/x),\qquad T=-W_{-1}(-\lambda x),
 \label{eq:Lambert}
\end{equation}
where $W_{-1}$ is the real lower branch of Lambert's function.  Equivalently,
$T>1$ is the large solution of $T\e^{-T}=\lambda x$.  Thus
\begin{equation}
 T=L+\log L-\log\lambda+O_\lambda(\log L/L).
 \label{eq:Texp}
\end{equation}
No complex branch convention is involved.

\begin{theorem}[Phase-corrected cap probability]
\label{thm:CDF}
Under \eqref{eq:BT}, uniformly in phase,
\begin{equation}
 t_x=T+O_{\at}(\log T/T).
 \label{eq:saddlelocation}
\end{equation}
Moreover,
\begin{align}
 \log C_\beta(x)
 &=-\frac{T^2}{2\lambda}
   +\left(\frac12+\frac1\lambda\right)T+R_\beta(T)
      +O_{\at}(\log^2 T),\label{eq:CDFT}\\
 &=-\frac{L^2}{2\lambda}-\frac{L\log L}{\lambda}
      +c_\lambda L+R_\beta(L)+O_{\at}(\log^2 L).
 \label{eq:CDFL}
\end{align}
\end{theorem}
\begin{proof}
At the saddle, Lemma~\ref{lem:tilted} gives
$x=\e^{-t_x}(t_x/\lambda)(1+O_{\at}(\log t_x/t_x))$.
Taking logarithms yields
\[
 t_x-\log t_x=L-\log\lambda+O_{\at}(\log t_x/t_x).
\]
First this implies $t_x\sim L$ uniformly.  Since
$T-\log T=L-\log\lambda$ and the derivative of $u-\log u$ is bounded below
by $1/2$ for $u\ge2$, it gives \eqref{eq:saddlelocation}.
Equation~\eqref{eq:Texp} follows either by one substitution in the same
implicit equation or by the identical derivative comparison.

Taking logarithms of the relative formula \eqref{eq:saddle} and inserting
Theorem~\ref{thm:profile} and Lemma~\ref{lem:tilted} gives
\[
 \log C_\beta(x)
 =-\frac{t_x^2}{2\lambda}
   +\left(\frac12+\frac1\lambda\right)t_x+R_\beta(t_x)
     +O_{\at}(\log^2 t_x).
\]
The variance prefactor contributes only $O(\log t_x)$ here.
Changing $t_x$ to $T$ changes the quadratic expression by $O_{\at}(\log T)$;
\eqref{eq:Lipschitz} bounds the change in the profile.  This proves
\eqref{eq:CDFT}.  Expand $T$ using \eqref{eq:Texp}, and note that
$R_\beta(T)-R_\beta(L)=O_\lambda(\log L)$.
The coefficient of $L$ is exactly the constant in \eqref{eq:clambda}, giving
\eqref{eq:CDFL}.
\end{proof}

\begin{theorem}[Uniform multiplicative inverse-quantile law]
\label{thm:quantile}
Assume \eqref{eq:BT}.  Put
\begin{equation}
 z=\sqrt{2\lambda\log(1/y)},\qquad
 Q_\lambda(y)=\frac z\lambda\exp(-1-\lambda/2-z).
 \label{eq:Qdef}
\end{equation}
Then, uniformly in $\beta$ as $y\downarrow0$,
\begin{equation}
 C_\beta^{-1}(y)=Q_\lambda(y)
     \exp\!\left(-\frac{\lambda R_\beta(z)}{z}\right)
     \left(1+O_{\at}\!\left(\frac{\log^2 z}{z}\right)\right).
 \label{eq:quantilemain}
\end{equation}
Equivalently,
\begin{equation}
 -\log C_\beta^{-1}(y)
 =z-\log z+1+\log\lambda+\frac\lambda2
       +\frac{\lambda R_\beta(z)}{z}
       +O_{\at}(\log^2 z/z).
 \label{eq:logquantile}
\end{equation}
\end{theorem}
\begin{proof}
Let $M=\log(1/y)$, and consider
\[
 E_\beta(L)=\frac{L^2}{2\lambda}+\frac{L\log L}{\lambda}
             -c_\lambda L-R_\beta(L).
\]
Equation~\eqref{eq:CDFL} says
$-\log C_\beta(\e^{-L})=E_\beta(L)+O_{\at}(\log^2 L)$.
Set
\[
 L_0=z-\log z+\lambda c_\lambda+\frac{\lambda R_\beta(z)}{z}.
\]
Since $R_\beta(z)\le z$, $L_0=z+O_\lambda(\log z)$ uniformly.
Expanding the elementary terms, and using the Lipschitz bound for $R_\beta$,
gives
$E_\beta(L_0)=M+O_{\at}(\log^2 z)$.
For $L$ near $L_0$, the elementary part of $E_\beta$ has slope at least
$z/(2\lambda)$ for sufficiently large $z$, while the profile has Lipschitz
constant $1$.  Hence increasing $L$ by $K\log^2 z/z$ increases $E_\beta(L)$
by more than the possible error when $K$ is sufficiently large; decreasing
$L$ has the reverse effect.  Monotonicity of $C_\beta$ brackets its inverse
between these two choices of $L$.
Thus $-\log C_\beta^{-1}(y)=L_0+O_{\at}(\log^2 z/z)$.
Insert $\lambda c_\lambda=1+\log\lambda+\lambda/2$ and exponentiate.
Since $\log^2 z/z\to0$, exponentiating the error gives the relative error in
\eqref{eq:quantilemain}.
\end{proof}

\begin{corollary}[A common asymptotic and a universal multiplicative range]
\label{cor:Qcommon}
Under \eqref{eq:BT}, condition \eqref{eq:polynomialphase} implies
\[
 C_\beta^{-1}(y)=Q_\lambda(y)
       \bigl(1+O_{\alpha,\lambda,c,\tau}(\log^2 z/z)\bigr).
\]
For arbitrary phases, the ratio $C_\beta^{-1}(y)/Q_\lambda(y)$ is
asymptotically confined to $[q,1]$, in the sense that its distance from this
interval tends to zero uniformly in $\beta$.
\end{corollary}
\begin{proof}
Apply Corollary~\ref{cor:polynomial} in \eqref{eq:quantilemain} for the first
statement, and use $0\le R_\beta(z)/z\le1$ for the second.
\end{proof}

\section{Phase classification: measure, category, and quantile oscillation}
\label{sec:phases}

For a nonexact phase define its inhomogeneous exponential approximation
exponent by
\begin{equation}
 \mu(\alpha,\beta)=\limsup_{n\to\infty}\frac{k_n(\beta)}n
                      \in[0,\infty].
 \label{eq:mu}
\end{equation}
By \eqref{eq:distance-amplitude}, replacing $k_n$ by
$(-\log d_n)_+$ leaves the exponent unchanged.

\begin{lemma}[Optimization of the profile]
\label{lem:profilespectrum}
For any irrational $\alpha$ and any nonexact phase,
\begin{equation}
 \limsup_{t\to\infty}\frac{R_\beta(t)}t
       =\frac{\mu(\alpha,\beta)}{\lambda+\mu(\alpha,\beta)},
 \label{eq:Rlimsup}
\end{equation}
where the right side is $1$ when $\mu=\infty$.  Under \eqref{eq:BT}, one
also has
\begin{equation}
 \liminf_{t\to\infty}\frac{R_\beta(t)}t=0.
 \label{eq:Rliminf}
\end{equation}
For an exact phase, instead, $R_\beta(t)/t\to1$.
\end{lemma}
\begin{proof}
For a fixed $n$ with finite $k_n$, the largest value of
$\min\{k_n,(t-\lambda n)_+\}/t$ is
$k_n/(\lambda n+k_n)$, attained at $t=\lambda n+k_n$.
For finite $\mu$, eventually $k_n\le(\mu+\eps)n$, giving the upper bound
$(\mu+\eps)/(\lambda+\mu+\eps)$ for all sufficiently late indices.
The finitely many earlier indices have bounded heights and contribute
$o(1)$ after division by $t$.  Indices with $k_n/n$ approaching its limsup
give the matching lower bound at their optimizing times.  The case
$\mu=\infty$ follows the same way from a sequence with $k_n/n\to\infty$;
the case $\mu=0$ follows already from the upper bound.

For \eqref{eq:Rliminf}, note that the sequence $k_n$ is unbounded by density
of the irrational orbit.  At a nonexact phase every term is finite, so it
has infinitely many strict record indices $m_j$.  Ignore initial records
with height zero.  Writing $d_j=d_{m_j}$, the next record has
$d_{j+1}<d_j$.  The two orbit points therefore satisfy
\[
 \frac{\kappa}{m_{j+1}-m_j}
 \le\dn{(m_{j+1}-m_j)\alpha}
 \le d_j+d_{j+1}
 \le2d_j\le\tfrac12\e^{-k_{m_j}},
\]
where the last inequality follows from $4d_j\le a_{m_j}=\e^{-k_{m_j}}$.
Hence $k_{m_j}\le\log m_{j+1}+C_\alpha$.
At the time $t_j=\lambda(m_{j+1}-1/2)$ the new record has not yet entered
the profile, so $R_\beta(t_j)\le k_{m_j}$ and
$R_\beta(t_j)/t_j\to0$.
Finally, if $k_m=\infty$ for an exact index, then
$R_\beta(t)\ge(t-\lambda m)_+$; combine this with $R_\beta(t)\le t$.
\end{proof}

\begin{theorem}[Endpoint and inverse-quantile spectrum]
\label{thm:spectrum}
Assume \eqref{eq:BT}, let $\beta$ be nonexact, and put
$p_\beta=\mu(\alpha,\beta)/(\lambda+\mu(\alpha,\beta))$ with the preceding
convention at infinity.  Then
\begin{align}
 \liminf_{t\to\infty}
   \frac{\log\LL_\beta(\e^t)+t^2/(2\lambda)}t&=\frac12,\nonumber\\
 \limsup_{t\to\infty}
   \frac{\log\LL_\beta(\e^t)+t^2/(2\lambda)}t&=\frac12+p_\beta,
 \label{eq:Lspec}\\
 \liminf_{L\to\infty}
   \frac{\log C_\beta(\e^{-L})+L^2/(2\lambda)+(L\log L)/\lambda}{L}
 &=c_\lambda,
 \label{eq:Cspecinf}\\
 \limsup_{L\to\infty}
   \frac{\log C_\beta(\e^{-L})+L^2/(2\lambda)+(L\log L)/\lambda}{L}
 &=c_\lambda+p_\beta,
 \label{eq:Cspecsup}\\
 \limsup_{y\downarrow0}\frac{C_\beta^{-1}(y)}{Q_\lambda(y)}&=1,\nonumber\\
 \liminf_{y\downarrow0}\frac{C_\beta^{-1}(y)}{Q_\lambda(y)}
 &=\exp(-\lambda p_\beta).
 \label{eq:Qspec}
\end{align}
In particular, $\mu=0$ gives $C_\beta^{-1}(y)\sim Q_\lambda(y)$, while
$\mu=\infty$ gives lower and upper quantile-ratio limits $q$ and $1$.
For an exact phase, the quantile ratio converges to $q$.
\end{theorem}
\begin{proof}
Divide \eqref{eq:laplaceprofile} and \eqref{eq:CDFL} by their large variables;
$\log^2 t/t\to0$.  Insert Lemma~\ref{lem:profilespectrum}.
Equation~\eqref{eq:quantilemain} and continuity and monotonicity of
$u\mapsto\exp(-\lambda u)$ give the quantile limits.
For an exact phase use the final conclusion of that lemma instead.
\end{proof}

\begin{theorem}[Almost every phase versus a thin residual set]
\label{thm:category}
Fix a badly approximable irrational $\alpha$ and $\lambda>0$.
For Lebesgue-almost every phase $\beta$, condition
\eqref{eq:polynomialphase} holds for every fixed $\tau>1$ with a suitable
constant, and
\begin{equation}
 C_\beta^{-1}(y)\sim Q_\lambda(y).
 \label{eq:a.e.Q}
\end{equation}
There is also a dense $G_\delta$ set of nonexact phases with $\mu=\infty$.
This set has Hausdorff dimension zero and for all its phases
\begin{equation}
 \liminf_{y\downarrow0}\frac{C_\beta^{-1}(y)}{Q_\lambda(y)}=q,
 \qquad
 \limsup_{y\downarrow0}\frac{C_\beta^{-1}(y)}{Q_\lambda(y)}=1.
 \label{eq:residualQ}
\end{equation}
Indeed, the entire set of nonexact phases with $\mu>0$ has Hausdorff
dimension zero.
\end{theorem}
\begin{proof}
For $\tau>1$, the phase event $d_n(\beta)<n^{-\tau}$ is an interval of length
at most $2n^{-\tau}$.  Its lengths are summable.  Borel--Cantelli implies
that almost every phase belongs to only finitely many of these events.
Remove the countable exact phases; the finitely many early positive
distances can be included by decreasing the constant $c$.
Taking a countable collection of rational exponents $\tau>1$ gives the
claim for every fixed real $\tau>1$ as well, by choosing a smaller rational
exponent.  Corollary~\ref{cor:Qcommon} proves \eqref{eq:a.e.Q}.

For integers $K,N\ge1$, let
\[
 U_{K,N}=\bigcup_{n\ge N}\{\beta\in\T:a_n(\beta)<\e^{-Kn}\}.
\]
Each is open and dense: it contains a neighborhood of every center
$1/2-n\alpha$ for $n\ge N$, and those centers are dense.
The intersection of all $U_{K,N}$ is a dense $G_\delta$ set on which
$\mu=\infty$.  Intersecting further with the complement of the countable
exact phases preserves the dense $G_\delta$ property.
Theorem~\ref{thm:spectrum} then proves \eqref{eq:residualQ}.

Finally, if $\mu>0$, then for some integer $K\ge1$ the inequality
$a_n<\e^{-n/K}$ holds infinitely often.  By
\eqref{eq:distance-amplitude}, each such event is contained in an interval
of length $O(\e^{-n/K})$.  For every $d>0$, the sum of the $d$th powers of
these lengths is finite.  Covering the limsup set by all intervals after
an arbitrary index $N$, and then letting $N\to\infty$, makes its
$d$-dimensional Hausdorff measure zero.  A countable union over $K$ still
has Hausdorff dimension zero.
\end{proof}

\begin{corollary}[A precise obstruction to a small carrier]
\label{cor:carrier}
For every fixed badly approximable $\alpha$, there are nonexact phases,
forming a dense $G_\delta$ set, such that
\[
 \limsup_{t\to\infty}
 \frac{\log\LL_\beta(\e^t)+t^2/(2\lambda)-t/2}{t}=1.
\]
Thus at these phases the remainder after the universal quadratic and linear
terms is neither bounded nor $O(\log t)$.
\end{corollary}
\begin{proof}
Use \eqref{eq:Lspec} with $\mu=\infty$ and subtract $1/2$.
\end{proof}

\begin{remark}[Interpretation of the source conjecture]
The preceding corollary does not contradict the source's displayed estimate;
that estimate holds in the stronger form of Corollary~\ref{cor:repo}.
It corrects the accompanying carrier expectation \emph{if} the word
nonresonant is taken to mean only that no factor is exactly zero.
Under a quantitative interpretation such as \eqref{eq:polynomialphase},
the exceptional profile is only logarithmic, and further control of the
remaining $O(\log^2 t)$ term is an open refinement rather than a contradiction.
\end{remark}

\section{Liouville rotations with superlinear corrections}
\label{sec:liouville}

We now drop \eqref{eq:BT} and fix the phase $\beta=1/2$.
Thus $a_n=2|\sin\pi n\alpha|$.  To isolate a genuine resonance effect rather
than the automatic linear term, define the baseline-corrected defect
\begin{equation}
 \Delta_\alpha(t)=B_\lambda(t)-H_{\alpha,1/2}(t).
 \label{eq:Delta}
\end{equation}
The extra $\alpha$ subscript makes variation in the rotation explicit.

\begin{lemma}[Uniform rational shadow and continuity bound]
\label{lem:rationalshadow}
For a reduced rational $p/r$, $r\ge2$,
\begin{equation}
 \Delta_{p/r}(t)=B_{\lambda r}(t)
      +O_\lambda((t+1)\log(2r))
 =\frac{t^2}{2\lambda r}
      +O_\lambda((t+1)\log(2r)+r).
 \label{eq:rationalshadow}
\end{equation}
The constants are independent of $p$ and $r$.  For any real $\alpha,\alpha'$,
\begin{equation}
 |H_{\alpha,1/2}(t)-H_{\alpha',1/2}(t)|
 \le K_\lambda\e^t|\alpha-\alpha'|,
 \qquad K_\lambda=\frac{\pi q}{(1-q)^2}.
 \label{eq:continuityalpha}
\end{equation}
\end{lemma}
\begin{proof}
At $\alpha=p/r$, exactly the multiples of $r$ vanish.  For all other
indices, multiplication by $p$ permutes the nonzero residues, so
\[
 2\sin(\pi/r)\le a_n\le2,
 \qquad 4/r\le a_n\le2.
\]
For $n\le\lfloor t/\lambda\rfloor$, replacing the nonzero factor $a_n$ by
$1$ changes its $h$ value by at most $|\log a_n|\le C\log(2r)$ because
$0\le h'\le1$.  The sum of all terms past that index is $O_\lambda(1)$
both before and after the replacement.  Consequently,
\[
 H_{p/r,1/2}(t)=
   \sum_{n\ge1}h(t-\lambda n)
  -\sum_{j\ge1}h(t-\lambda rj)
  +O_\lambda((t+1)\log(2r)).
\]
Use \eqref{eq:scalarcut}, which is uniform for lattice spacings
$\lambda r\ge\lambda$, followed by \eqref{eq:Bexact}.

For the continuity bound,
\[
 |2q^n|\sin\pi n\alpha|-2q^n|\sin\pi n\alpha'||
 \le2\pi nq^n|\alpha-\alpha'|.
\]
Since $F$ is $1/2$-Lipschitz by Lemma~\ref{lem:kernel}, sum the resulting
bounds for $F(\e^t b_n)$ and use
$\sum_{n\ge1}nq^n=q/(1-q)^2$.
\end{proof}

\begin{theorem}[Dense Liouville families with arbitrarily near-quadratic spikes]
\label{thm:liouville}
Fix $\lambda>0$.  There is a dense $G_\delta$ set of irrational Liouville
numbers $\alpha$ such that the following holds simultaneously for every
integer $k\ge2$: there is a sequence $t_j\to\infty$ with
\begin{equation}
 \Delta_\alpha(t_j)
 \sim\frac{1}{2\lambda}\,t_j^{\,2-1/k}.
 \label{eq:liouvillespikes}
\end{equation}
The phase can be fixed at $\beta=1/2$ for the whole construction.
In particular, the correction is larger than every fixed multiple of $t$
along the $k=2$ subsequence, even after subtracting the exact scalar baseline.
\end{theorem}
\begin{proof}
For integers $k\ge2$ and $N\ge1$, define
\[
 \mathcal U_{k,N}=
 \bigcup_{\substack{r\ge N,\ p\in\mathbb Z\\(p,r)=1}}
 \left(\frac pr-\e^{-2r^k},\frac pr+\e^{-2r^k}\right).
\]
These sets are open and dense in $\R$, since reduced rationals with
arbitrarily large denominators are dense.  Their countable intersection
$\mathcal A=\bigcap_{k\ge2}\bigcap_{N\ge1}\mathcal U_{k,N}$ is dense
$G_\delta$.  A rational $a/b$ cannot belong: for sufficiently large $r$,
its distance to a different reduced $p/r$ is at least $1/(br)$, which
exceeds $\e^{-2r^k}$.  Every member is therefore irrational and admits
approximations stronger than $r^{-M}$ for every $M$, hence is Liouville.

Fix $\alpha\in\mathcal A$ and an integer $k\ge2$.
Choose reduced $p_j/r_j$ with $r_j\to\infty$ and
$|\alpha-p_j/r_j|<\e^{-2r_j^k}$, and set $t_j=r_j^k$.
By \eqref{eq:continuityalpha},
\[
 |\Delta_\alpha(t_j)-\Delta_{p_j/r_j}(t_j)|
 \le K_\lambda\e^{-t_j}.
\]
The rational estimate gives
\[
 \Delta_{p_j/r_j}(t_j)
 =\frac{r_j^{2k-1}}{2\lambda}
   +O_\lambda(r_j^k\log(2r_j)+r_j).
\]
The error is $o(r_j^{2k-1})$ because $k\ge2$.
Since $r_j^{2k-1}=t_j^{2-1/k}$, the desired asymptotic follows.
The definition of $\mathcal A$ ensures that this works for every $k$,
with different subsequences allowed.
\end{proof}

\begin{remark}[Consistency with irrational universality]
For each fixed $k$, the exponent $2-1/k$ is strictly below $2$.
Thus these large corrections are compatible with
$H_\alpha(t)\sim t^2/(2\lambda)$ from Theorem~\ref{thm:universal}.
The effective rational fraction of missing digits is $1/r_j$, which tends
to zero, although its integrated contribution $t_j^2/(2\lambda r_j)$
can be much larger than a linear term.
\end{remark}

\begin{remark}[What is not transferred to Liouville cap quantiles]
The relative saddle proof above uses a bounded-type separation estimate to
produce a gamma block with only logarithmically many exceptions.  It has not
been extended here to the Liouville family.  The Liouville result concerns
Laplace corrections; the leading cap and inverse laws still follow from
Theorem~\ref{thm:universal}, but a sharp Liouville probability or quantile
spike formula is not asserted.
\end{remark}

\section{Reproducible calculations and validation limits}
\label{sec:computations}

The accompanying program \path{verify_endpoints.py} evaluates the elementary
infinite sums directly in logarithmic coordinates.  It requires Python and
\texttt{mpmath}, uses no network access or random sampling, and writes its
results to CSV and JSON.  The mathematical proofs do not depend on these
calculations.  The tables below are finite examples, not evidence sufficient
to infer an infinite-subsequence theorem.

\subsection{Bounded-type and deliberately resonant phases}

Take $q=1/2$, $\lambda=\log2$, and
$\alpha=(\sqrt5-1)/2$.  Alongside $\beta=0$, consider the phase
\begin{equation}
 \beta_* =\frac12-20\alpha+\e^{-120}\pmod1.
 \label{eq:engineeredphase}
\end{equation}
Its twentieth amplitude is
$2\sin(\pi\e^{-120})$, so
$k_{20}=120-\log(2\pi)+O(\e^{-240})$.
Before the corresponding height is reached, its contribution to the profile
is $t-20\log2$; after that it saturates.  This is a controlled demonstration
of one long ramp, not a construction of infinitely many ramps at a fixed
phase.  The latter is supplied by the Baire argument, not by the example.

\begin{center}\small
\begin{tabular}{@{}llrrr@{}}
\toprule
Phase & $t$ & $R_\beta(t)$ & $H-B+R$ & $H^\prime/(t/\lambda)$ \\
\midrule
$\beta=0$ & 40 & 3.186248 & 5.003670 & 0.994760 \\
$\beta=0$ & 80 & 4.629816 & 5.743629 & 0.995093 \\
$\beta=0$ & 120 & 4.629816 & 5.770088 & 0.998650 \\
$\beta=0$ & 160 & 4.629816 & 6.306261 & 0.997050 \\
engineered phase & 40 & 26.137056 & 5.767516 & 0.969153 \\
engineered phase & 80 & 66.137056 & 5.824267 & 0.986599 \\
engineered phase & 120 & 106.137056 & 6.067652 & 0.991524 \\
engineered phase & 160 & 118.162123 & 7.265925 & 0.998964 \\
\bottomrule
\end{tabular}

\end{center}

The residual $H-B+R$ is modest at the displayed scales, while omitting $R$
would miss a correction exceeding $100$ in the deliberately resonant case.
The table does not establish an optimal residual order or numerical constants
for the asymptotic theorems.  The second ordinary phase $\beta=1/2$ and
additional values are retained in the CSV file.

\subsection{Rational shadows for the Liouville construction}

For $\alpha=1/r$, $\beta=1/2$, and $t=r^2$, the theorem predicts
$\Delta_{1/r}(t)\sim t^2/(2\lambda r)$ as $r\to\infty$.
Multiples of $r$ are set to zero exactly in the code rather than by rounded
sine evaluation.

\begin{center}\small
\begin{tabular}{@{}rrrrr@{}}
\toprule
$r$ & $t=r^2$ & $\Delta_{1/r}(t)$ & $t^2/(2\lambda r)$ & ratio \\
\midrule
7 & 49 & 202.642024 & 247.422200 & 0.819013 \\
13 & 169 & 1450.989793 & 1584.800502 & 0.915566 \\
23 & 529 & 8410.452487 & 8776.635281 & 0.958278 \\
41 & 1681 & 48661.591192 & 49715.992457 & 0.978792 \\
\bottomrule
\end{tabular}

\end{center}

For any irrational $\alpha'$ satisfying
$|\alpha'-1/r|<\e^{-2r^2}$, Lemma~\ref{lem:rationalshadow} transfers the
corresponding value of $H$ with error at most $K_\lambda\e^{-r^2}$.
This continuity bound is analytic.  The table alone does not exhibit one
irrational number having all these approximations; the nested dense-open
construction in Theorem~\ref{thm:liouville} supplies suitable single numbers.

\subsection{Kernel, precision, and saddle checks}

The program checks the finite-sum and fractional-part versions of
\eqref{eq:Bexact}; numerically differentiates the kernel to check
$h'=J$ and $J-uJ'=V$; tests one profile increment against its monotonicity and
Lipschitz bounds; and recomputes the deliberately resonant example at higher
precision.  With the default $110$ decimal digits and an additional
$35$ digits for that comparison, the absolute difference in $H_{\beta_*}(120)$
was approximately $5.51\times10^{-65}$.

For $\beta=0$ and $x=\e^{-20}$, deterministic bisection gave
\[
 T=23.52455780924499\ldots,\qquad
 t_x=23.51292579548681\ldots,
\]
with a relative residual in the saddle equation below $6\times10^{-31}$.
The saddle expression for $\log C_0(\e^{-20})$ evaluates to approximately
$-357.5854846402$.  This last number is explicitly a \emph{saddle
approximation}: the program does not independently evaluate the exact cap
CDF, so it supplies no finite-$x$ relative error certificate for that number.

The analytic omitted-tail bounds in the main tables are below $10^{-40}$.
Rounding error is separate: calculations use high precision and a precision
comparison, not interval arithmetic.  The JSON report records these limits,
the actual equation residual, and the checks performed.  No Lean, Coq, or
other kernel verification is claimed.

\subsection{Dependency and formalization roadmap}

The proof dependencies are short enough to separate into four layers.
The distributional layer consists of the uniform-series construction,
Laplace product, tilted summands, and scalar kernel inequalities.  The
arithmetic layer consists of the bounded-type spacing and interval-discrepancy
bounds, followed by the trimmed singular sum.  The analytic layer passes
from that sum to the profile expansion and the gamma-block local limit.
The final layer is real-variable inversion and the measure/category
classification.  The Liouville argument branches off after the kernel layer
and uses rational comparison, continuity, and Baire category instead.

A practical formalization should first establish the deterministic finite
profile and the kernel bounds, which also support certified computation.
Next it should formalize the continued-fraction block discrepancy and the
truncation integral.  Only then is the infinite-product and probability
infrastructure needed.  The relative local-limit step should be formalized
as the standalone gamma-convolution lemma rather than hidden inside an
unproved appeal to a saddle-point method.  A separate classical-analysis
proof of monotone inverse bracketing suffices for the quantile theorem.
These are proposed proof interfaces, not untested Lean code presented as
working implementation.

\section{Further research questions and topics}
\label{sec:future}

The questions below distinguish refinements of this article from broader
problems requiring an additional literature check.  In particular, a question
about importing a sharper shifted Sudler estimate is not a claim that the
corresponding estimate is unknown in all existing work.

\begin{problem}[Sharp uniform residual after extracting the profile]
Can the $O_{\at}(\log^2 t)$ residual in
\eqref{eq:mainprofile} be replaced by $O_{\at}(\log t)$, or by a smaller
explicit arithmetic term plus a bounded remainder?  First determine which
available shifted Sudler-product estimates apply with constants uniform in
phase after the nearest singular term has been removed.  The target is the
cap expansion, not a rediscovery of an existing unshifted product estimate.
\end{problem}

\begin{problem}[Renormalized carriers at quadratic rotations]
For a quadratic irrational $\alpha$ and an explicitly specified arithmetic
class of phases, identify the behavior of
$H_\beta(t)-B_\lambda(t)+R_\beta(t)$ along continued-fraction scales.
Is the natural carrier periodic in a renormalization coordinate,
quasiperiodic, or better described by an iterated functional relation?
Any proposed bounded carrier should state the phase class and the variable
in which boundedness is claimed.
\end{problem}

\begin{problem}[Effective finite-exponent phase realization]
For a fixed badly approximable $\alpha$ and prescribed
$0<\mu<\infty$, construct or import a construction of nonexact phases with
$\mu(\alpha,\beta)=\mu$, while controlling the optimizing subsequences.
Theorem~\ref{thm:spectrum} then gives the corresponding quantile-ratio
interval explicitly.  An effective construction with error bounds would add
more information than the conditional classification proved here.
\end{problem}

\begin{problem}[Sharp Liouville cap and inverse spikes]
Extend the rational-shadow argument beyond the Laplace transform to a
relative tilted local limit on the subsequences $t_j=r_j^k$.
Determine the resulting deviations of $C_{\alpha,1/2}(x)$ and its inverse
from their universal leading laws.  A promising route is to separate the
nearly vanished rational-residue terms from a large remaining gamma block;
the required estimates must be uniform in the increasing denominator.
\end{problem}

\begin{problem}[Several singular directions and matrix cap laws]
Replace one cosine amplitude by a finite family of projected rotating
vectors, or by a periodic amplitude with several zeros of specified orders.
Determine whether a sum of finitely many truncated resonance profiles gives
a phase-uniform expansion.  For genuine matrix dilations, investigate how
multiple rotation frequencies, unequal contraction rates, and polynomial
Jordan factors change both the quadratic coefficient and the
$L\log L$ correction.
\end{problem}

\begin{problem}[Quantitative local errors and joint parameter limits]
Derive explicit constants and rates in the gamma-block approximation and in
\eqref{eq:saddle}.  Then study a joint limit $x\downarrow0$ and $q\uparrow1$.
The present constants are allowed to depend on $\lambda$, so the theorems
cannot simply be used uniformly as $\lambda\downarrow0$.  Locating the
transition between the cap regime and the nearly Gaussian central regime
would connect endpoint and bulk asymptotics.
\end{problem}

\begin{problem}[Statistics under phase averaging]
For a uniformly distributed phase, estimate the distribution of
$R_\beta(t)$, its moments, and the tails of the quantile ratio at a finite
scale.  Almost-everywhere asymptotics and Hausdorff dimension zero do not
supply quantitative finite-scale probabilities.  The truncated nature of the
profile makes this a bounded observable even at exact phases.
\end{problem}

\begin{problem}[Certified computation and proof-assistant integration]
Implement interval-certified evaluation of $H_\beta$, $H_\beta'$, and
$\VV_\beta$, combined with a rigorous saddle bracket and explicit local-limit
bounds.  Exploit the finite maximum in \eqref{eq:profiledef} to avoid having
to decide whether a computable phase is exactly singular.  Formalize the
kernel, discrepancy, trimming, and inverse-bracketing interfaces in the
repository before attempting a monolithic formalization of the full saddle
theorem.
\end{problem}

\section{Conclusions}

The leading quadratic endpoint scale is much more robust than the singular
cocycle alone suggests: every irrational rotation has it, uniformly in phase.
Bounded type gives a stronger structural conclusion.  After a deterministic
scalar baseline is removed, a single truncated deepest-return profile controls
the Laplace transform up to a squared-logarithmic error.  This control is
strong enough, through a relative tilted local limit, to determine a
multiplicative inverse-quantile asymptotic.

The arithmetic distinction is not merely between rational, Diophantine, and
Liouville angles.  Even at one fixed badly approximable angle, the phase has
its own exponential approximation exponent.  It governs a real, explicitly
quantified oscillation in extreme quantiles.  Almost-everywhere universality
coexists with a dense, dimension-zero exceptional set.  Liouville angles
introduce a further mechanism: long rational shadows suppress whole residue
classes and produce near-quadratic corrections along selected subsequences.

Thus the two concrete repository questions lead to both positive theorems and
a necessary qualification.  The displayed Diophantine bound is valid in a
stronger form, and Liouville spikes exist in a stronger baseline-corrected
form.  A bounded or logarithmic carrier, however, needs a quantitative phase
hypothesis and cannot follow from absence of exact zeros alone.

\appendix
\section{Source provenance, normalization checks, and package contents}
\label{app:provenance}

The inspected repository is \url{https://github.com/VladimirReshetnikov/ProveIt}.
The relevant report is located at
\begin{quote}\small
\path{Analysis/FabiusFunction/docs/semi-formalized-research-frontiers/drafts/representations/Matrix_Dilated_Fabius_Rvachev_Frontier_Report/matrix_dilated_fabius_rvachev.tex}.
\end{quote}
The pinned repository revision is \pinned, and the report blob is
\texttt{143f6da778dfc0d5ab04133005ba6715a3cf155b}.
The defining cap product appears in the section on directional caps; the two
question labels are \texttt{conj:diophantine-endpoint} and
\texttt{conj:liouville}, inspected in source lines approximately
1565--1610 of this revision.  The report is a repository research document,
not an independently refereed source of the new theorems in this article.

The present normalization uses $n\ge1$, the factor
$2|\cos\pi(n\alpha+\beta)|$, and the convention $q=\e^{-\lambda}$.
The logarithmic phase average is zero:
$\int_\T\log(2|\cos\pi u|)\dd u=0$.
Consequently the common quantile $Q_\lambda$ corresponds to the scalar
baseline with weights $q^n$, not to a silently rescaled scalar law with
weights $2q^n$.  A deterministic external scale multiplies every quantile
by that scale, as in Proposition~\ref{prop:cap}.  Exact zero factors are
omitted from products, assigned $k_n=\infty$ in profiles, and are handled
separately from nonexact phases in the exponent classification.

The archive contains the self-contained LaTeX driver \path{article.tex}, its
compiled \path{article.pdf}, the deterministic verification program,
\path{requirements.txt}, a README with reproduction commands, the two CSV
result files, their LaTeX table counterparts, a JSON verification report,
and a machine-readable provenance record.  No external repository source,
font file, or unverified proof-assistant script is needed to build the article.
The only LaTeX inclusions are the two small table files supplied in
\path{data/}.

\clearpage
\phantomsection
\addcontentsline{toc}{section}{References}
\begin{thebibliography}{9}

\bibitem{RepoMatrix}
\repo\ repository research report.
\emph{Matrix-Dilated Fabius--Rvachev Laws: Infinite Box Splines,
Thue--Morse Directional Derivatives, Bell--Bernoulli Tensor Resolvents,
and Rotating-Zonoid $q$-Series}.
Report dated 30 August 2026; inspected at the repository revision recorded in
Appendix~\ref{app:provenance}.  In particular, the directional-cap section
and source labels \texttt{conj:diophantine-endpoint} and
\texttt{conj:liouville}.
\url{https://github.com/VladimirReshetnikov/ProveIt/blob/fbba58593dc0622aa914972896150d4848f935b5/Analysis/FabiusFunction/docs/semi-formalized-research-frontiers/drafts/representations/Matrix_Dilated_Fabius_Rvachev_Frontier_Report/matrix_dilated_fabius_rvachev.tex}.

\bibitem{Arias}
J.~Arias de Reyna.
\emph{An infinitely differentiable function with compact support:
Definition and properties}.
ArXiv:1702.05442, 2017 posting of the author's earlier article.
\url{https://arxiv.org/abs/1702.05442}.

\bibitem{VerschuerenMestel}
P.~Verschueren and B.~Mestel.
\emph{On the growth of Sudler's sine product at the golden rotation number}.
ArXiv:1411.2252, 2014.
\url{https://arxiv.org/abs/1411.2252}.

\bibitem{Borda}
B.~Borda.
\emph{On the distribution of Sudler products and Birkhoff sums for the
irrational rotation}.
ArXiv:2104.06716, version 2, 2023.
\url{https://arxiv.org/abs/2104.06716}.

\bibitem{Hauke}
M.~Hauke.
\emph{On the asymptotic behaviour of Sudler products for badly
approximable numbers}.
ArXiv:2210.07743, 2022.
\url{https://arxiv.org/abs/2210.07743}.

\end{thebibliography}
\end{document}
